\section{A Program-Level Rule for AST}
\label{sec:rule}

Consider $C=\while{b}{B}$ with loop-free body $B$ and a set $M$ of
initial subdistributions.

\subsection{The Proof Rule}

\begin{theorem}[Program-level AST rule]
\label{thm:program-rule}
The loop $C$ is AST on $M$ if a distributional invariant $I$ admits
functions $v:S\to\mathbb R_{\geq0}$ and $u:S\to\mathbb N$, where
$S=\supp(I)$, with the following properties:
\begin{enumerate}
 \item $v(\sigma)=0$ exactly when $\sigma\models\neg b$, and
       $u(\sigma)=0$ whenever $\sigma\models\neg b$;
 \item for every $\sigma\in S\cap\llbracket b\rrbracket$,
 \[
   \sum_{\sigma'}\post{B}(\dirac{\sigma})(\sigma')v(\sigma')
     \leq v(\sigma);
 \]
 \item some $\varepsilon>0$ satisfies, for every
 $\sigma\in S\cap\llbracket b\rrbracket$,
 \[
   \sum_{\sigma':\,u(\sigma')<u(\sigma)}
      \post{B}(\dirac{\sigma})(\sigma')\geq\varepsilon;
   \tag{Decrease}
 \]
 \item for every $r\geq0$, the values of $u$ on $v_{\leq r}$ are bounded.
\end{enumerate}
\end{theorem}
The decrease condition itself forces $u(\sigma)>0$ on enabled states.
It permits increases on other outcomes, including rejection resets.
Sublevel boundedness, rather than global boundedness, admits
null-recurrent computations.

\paragraph{Why the rule works.}
Fix $v_{\leq r}$ and a positive integer bound $K_r$ on $u$ there.
Starting in this set, the run has probability at least
$\min(\varepsilon,1)^{K_r}$ of either terminating or leaving the set
within $K_r$ iterations: otherwise $K_r$ successive decreases reach
$u=0$, which cannot be enabled. Repeating this bound rules out
nontermination confined to the set. Nonnegative supermartingale
convergence makes unbounded values of $v$ a probability-zero event.
The formal proof in \Cref{sec:metatheory} establishes the same result
by constructing the pCFG certificates.

\subsection{Illustration: The Symmetric Random Walk}
\label{sec:walk}

Consider the symmetric random walk
\[
 \ass{x}{k};\quad
 \while{x>0}{\pchoice{\ass{x}{x-1}}{1/2}{\ass{x}{x+1}}},
 \qquad k\in\mathbb{N}_{>0}.
\]
This loop terminates almost surely at $x=0$, although its expected termination time is infinite. It is therefore a compact illustration of why the supermartingale and variant play different roles.

At loop entry, let $M=\{\dirac{\sigma}\mid\sigma(x)\in\mathbb N_{>0}\}$ and
$I=\{\mu\in\subdists\mid\supp(\mu)\subseteq\{\sigma\mid\sigma(x)\in\mathbb{N}\}\}$.
The set $I$ is distributionally inductive. Choose
\[
 v(\sigma)=u(\sigma)=\sigma(x).
\]
For every enabled state,
\[
 \sum_{\sigma'}\post{B}(\dirac{\sigma})(\sigma')v(\sigma')
 =\tfrac12(x-1)+\tfrac12(x+1)=x=v(\sigma),
\]
so $v$ is in fact a martingale. With probability $1/2$, the left branch is taken and $u$ decreases by one; hence the decrease obligation holds with $\varepsilon=1/2$. Finally, on $v_{\le r}$ we have $u\le r$. \Cref{thm:program-rule} proves AST.

At pCFG level, the same argument requires location-dependent values for the guard, probabilistic choice, and assignments. The program-level rule expresses it using the two state functions above and one body-step expectation. We make this representational comparison explicit for FDR in \Cref{sec:fdr}.

\Cref{sec:synthesis-examples} returns to this example and synthesizes the two functions automatically within the chosen template space.

\subsection{Iteration-Level Reasoning and Modularity}

The difference from \Cref{thm:pcfg-rule} is the granularity of progress:
our obligations inspect the final distribution of a body execution.
Internal assignments, tests, and choices are composed through
$\post{B}$, so one need not supply a certificate at each intermediate
location. Distributional invariants can be reused, and simplification of
$\post{B}(\dirac{\sigma})$ exposes the outcomes relevant to descent.

The pCFG rule permits a different progress bound $\varepsilon_r$ on each
sublevel set. For our fixed choices and loop-free bodies, every feasible
body path has a uniform positive probability lower bound
(\Cref{eq:body-path-bound}); hence one decreasing path per state suffices
for a global $\varepsilon$. The sampler proofs below use $\varepsilon=1/2$.

Terminating preprocessing is handled by transforming $M$ into the
loop-entry distributions. A finite data structure produced beforehand
is a read-only parameter of the sampling loop; its construction is
verified separately by ordinary termination reasoning. This isolates
the AST argument to the iteration that performs the sampling.
